\lecture{30}{Алгоритмы консенсуса}

\sourcevideo
    {https://www.youtube.com/playlist?list=PLOzRYVm0a65fhKDS8cYIC7YCuVGJ4FOJx}
    {Week 9: Lecture 30: Consensus Algorithms}

В дискретном среднем консенсусе требуется построить матрицу смешивания,
которая сохраняет сумму состояний и подавляет рассогласование. Для
статических неориентированных графов это можно сделать локально с
помощью весов Метрополиса. В непрерывном времени ту же структуру
задаёт лапласовская динамика
\[
    \dot x=-Lx.
\]

\section{Дважды стохастические веса}

Пусть
\[
    x^{k+1}=Wx^k,
    \qquad
    W=(w_{ij})_{i,j=1}^{N}\ge0.
\]
Строчная и столбцовая стохастичность записываются как
\[
    W\mathbf1=\mathbf1,
    \qquad
    \mathbf1^{\mathsf T}W=\mathbf1^{\mathsf T}.
\]
При выполнении второго равенства сумма состояний сохраняется:
\[
    \mathbf1^{\mathsf T}x^{k+1}
    =
    \mathbf1^{\mathsf T}x^k.
\]
Если одновременно достигается консенсус
\(x^k\to c\mathbf1\), то
\[
    c=\frac1N\sum_{i=1}^{N}x_i^0.
\]
Следовательно, двойная стохастичность превращает обычный консенсус в
средний. Для сходимости дополнительно требуется апериодичность:
матрица
\[
    \begin{pmatrix}
        0&1\\
        1&0
    \end{pmatrix}
\]
дважды стохастична и имеет связный граф, но порождает периодические
колебания. На связном неориентированном графе апериодичность обычно
обеспечивается положительными диагональными весами.

Один общий способ построения дважды стохастической матрицы ---
масштабирование Синкхорна--Кноппа. Для неотрицательной матрицы
\(M^0\) последовательно нормируют строки и столбцы:
\[
    r_i^k=\sum_jM_{ij}^k,
    \qquad
    \widetilde M_{ij}^k=\frac{M_{ij}^k}{r_i^k},
\]
\[
    c_j^k=\sum_i\widetilde M_{ij}^k,
    \qquad
    M_{ij}^{k+1}=\frac{\widetilde M_{ij}^k}{c_j^k}.
\]
В матричной форме это диагональное масштабирование
\[
    M^k\longmapsto R_kM^k\longmapsto R_kM^kC_k.
\]
Сходимость справедлива не для любой неотрицательной матрицы: носитель
должен допускать дважды стохастическое масштабирование; стандартное
достаточное условие формулируется через полный носитель, а полная
неразложимость обеспечивает единственность масштабирования с точностью
до общего множителя.

Алгоритм требует повторного вычисления нормировок по всей матрице и
поэтому не является непосредственно локальным. Он также не решает
дискретную задачу назначения: матрицы перестановок являются частным
случаем дважды стохастических матриц, но масштабирование лишь переводит
положительные оценки в соответствующее непрерывное множество.

\section{Веса Метрополиса}

Пусть \(\mathcal G=(\mathcal V,\mathcal E)\) --- связный
неориентированный граф, а
\[
    d_i=|\mathcal N_i|
\]
--- степень вершины \(i\). Для ребра \(\{i,j\}\in\mathcal E\)
положим
\[
    w_{ij}
    =
    \frac1{1+\max\{d_i,d_j\}},
\]
для несмежных различных вершин --- \(w_{ij}=0\), а диагональный вес
определим формулой
\[
    w_{ii}
    =
    1-\sum_{j\in\mathcal N_i}w_{ij}.
\]
Для вычисления этих коэффициентов агенту нужны только собственная
степень и степени соседей, которые можно передать за один
предварительный раунд.

По симметрии выражения
\[
    w_{ij}=w_{ji},
\]
поэтому \(W=W^{\mathsf T}\). По определению диагонального веса
\[
    W\mathbf1=\mathbf1,
\]
а симметричность даёт
\[
    \mathbf1^{\mathsf T}W=\mathbf1^{\mathsf T}.
\]
Кроме того,
\[
    w_{ij}
    \le
    \frac1{1+d_i},
    \qquad
    \sum_{j\in\mathcal N_i}w_{ij}
    \le
    \frac{d_i}{1+d_i}<1,
\]
следовательно, \(w_{ii}>0\). Поэтому на связном графе матрица
примитивна и
\[
    W^k
    \longrightarrow
    \frac1N\mathbf1\mathbf1^{\mathsf T}.
\]

Например, для графа с рёбрами
\[
    \mathcal E
    =
    \{\{1,2\},\{2,3\},\{2,4\},\{3,4\}\}
\]
имеем
\[
    (d_1,d_2,d_3,d_4)=(1,3,2,2)
\]
и
\[
    W=
    \begin{pmatrix}
        \frac34&\frac14&0&0\\[1mm]
        \frac14&\frac14&\frac14&\frac14\\[1mm]
        0&\frac14&\frac5{12}&\frac13\\[1mm]
        0&\frac14&\frac13&\frac5{12}
    \end{pmatrix}.
\]
Локальная итерация принимает вид
\[
    x_i^{k+1}
    =
    w_{ii}x_i^k
    +
    \sum_{j\in\mathcal N_i}w_{ij}x_j^k,
\]
то есть использует только собственное состояние и состояния соседей.

\section{Непрерывный лапласовский консенсус}

Пусть \(A=(a_{ij})\) --- симметричная матрица неотрицательных весов
неориентированного графа,
\[
    D=\operatorname{diag}(d_1,\ldots,d_N),
    \qquad
    d_i=\sum_ja_{ij},
\]
а
\[
    L=D-A
\]
--- матрица Лапласа. Непрерывный алгоритм задаётся системой
\[
    \dot x=-Lx.
\]
Её локальная форма равна
\[
    \dot x_i
    =
    \sum_{j\in\mathcal N_i}a_{ij}(x_j-x_i).
\]
Таким образом, матричная запись нужна для анализа, а реализация
остаётся полностью локальной.

Для пути из трёх вершин
\[
    A=
    \begin{pmatrix}
        0&1&0\\
        1&0&1\\
        0&1&0
    \end{pmatrix},
    \qquad
    D=
    \begin{pmatrix}
        1&0&0\\
        0&2&0\\
        0&0&1
    \end{pmatrix},
\]
поэтому
\[
    L=
    \begin{pmatrix}
        1&-1&0\\
        -1&2&-1\\
        0&-1&1
    \end{pmatrix}
\]
и
\[
    \dot x_1=x_2-x_1,
    \qquad
    \dot x_2=(x_1-x_2)+(x_3-x_2),
    \qquad
    \dot x_3=x_2-x_3.
\]

Для связного неориентированного графа
\[
    \ker L=\operatorname{span}\{\mathbf1\},
\]
поэтому равновесия имеют вид \(x=c\mathbf1\). Поскольку
\[
    \mathbf1^{\mathsf T}L=0,
\]
сумма и среднее состояний сохраняются:
\[
    \sum_{i=1}^{N}x_i(t)
    =
    \sum_{i=1}^{N}x_i(0),
    \qquad
    \bar x(t)=\bar x(0).
\]
Следовательно, любой достигнутый консенсус является средним.

\section{Сходимость и её скорость}

Решение линейной системы имеет вид
\[
    x(t)=\exp(-Lt)x(0).
\]
Для связного графа спектр симметричного Лапласиана удовлетворяет
\[
    0=\lambda_1(L)
    <
    \lambda_2(L)
    \le\cdots\le
    \lambda_N(L),
\]
а нормированный собственный вектор для нулевого собственного значения
равен
\[
    q_1=\frac1{\sqrt N}\mathbf1.
\]
Из спектрального разложения
\[
    \exp(-Lt)
    =
    \sum_{i=1}^{N}
    e^{-\lambda_i(L)t}q_iq_i^{\mathsf T}
\]
получаем
\[
    \lim_{t\to\infty}\exp(-Lt)
    =
    \frac1N\mathbf1\mathbf1^{\mathsf T}
\]
и, следовательно,
\[
    x(t)
    \longrightarrow
    \frac1N\mathbf1\mathbf1^{\mathsf T}x(0).
\]

Положим
\[
    e(t)=x(t)-\bar x\mathbf1,
    \qquad
    V(e)=\frac12\lVert e\rVert^2.
\]
Так как \(e\perp\mathbf1\), то
\[
    \dot V
    =
    -e^{\mathsf T}Le
    \le
    -\lambda_2(L)\lVert e\rVert^2
    =
    -2\lambda_2(L)V.
\]
Отсюда
\[
    V(t)
    \le
    e^{-2\lambda_2(L)t}V(0)
\]
и
\[
    \lVert x(t)-\bar x\mathbf1\rVert
    \le
    e^{-\lambda_2(L)t}
    \lVert x(0)-\bar x\mathbf1\rVert.
\]
Скорость непрерывного консенсуса определяется алгебраической
связностью \(\lambda_2(L)\).

Знак минус в динамике существенен. При \(\dot x=Lx\) компоненты в
собственных направлениях с \(\lambda_i(L)>0\) растут как
\(e^{\lambda_i(L)t}\), поэтому рассогласование увеличивается.

\section{Дискретизация и векторные состояния}

Явная схема Эйлера с шагом \(h>0\) даёт
\[
    x^{k+1}
    =
    x^k-hLx^k
    =
    (I-hL)x^k.
\]
Следовательно,
\[
    W=I-hL.
\]
Матрица \(W\) симметрична и сохраняет среднее. Локально
\[
    x_i^{k+1}
    =
    (1-hd_i)x_i^k
    +
    h\sum_{j\in\mathcal N_i}a_{ij}x_j^k.
\]
Для неотрицательности весов достаточно
\[
    0<h\le\frac1{d_{\max}},
    \qquad
    d_{\max}=\max_i d_i,
\]
а спектральная устойчивость требует
\[
    0<h<\frac2{\lambda_N(L)}.
\]
Первое условие обеспечивает стохастическую интерпретацию, второе ---
сжатие всех неконсенсусных мод; при стандартных невзвешенных графах
первое условие автоматически является безопасным и для устойчивости.

Если агент хранит вектор \(x_i(t)\in\mathbb R^d\), то после
объединения состояний
\[
    \mathbf x(t)
    =
    (x_1(t)^{\mathsf T},\ldots,x_N(t)^{\mathsf T})^{\mathsf T}
\]
динамика записывается как
\[
    \dot{\mathbf x}
    =
    -(L\otimes I_d)\mathbf x.
\]
Для связного графа
\[
    \mathbf x(t)
    \longrightarrow
    \mathbf1\otimes
    \left(
        \frac1N\sum_{i=1}^{N}x_i(0)
    \right),
\]
то есть каждая координата усредняется независимо.

\section{Реализация и границы модели}

Для локального вычисления
\[
    \dot x_i
    =
    \sum_{j\in\mathcal N_i}a_{ij}(x_j-x_i)
\]
агенту нужны только собственное состояние, состояния соседей и веса
инцидентных рёбер. Полная матрица Лапласа и глобальная топология ему
не требуются.

Полученные гарантии относятся к статическому связному
неориентированному графу с симметричными неотрицательными весами. В
ориентированной сети лапласиан обычно несимметричен и среднее может не
сохраняться. Для изменяющихся графов требуется совместная связность на
временных окнах, а задержки, потери сообщений и асинхронность требуют
отдельного анализа. Линейная динамика достигает консенсуса лишь
асимптотически; конечно-временные и фиксированно-временные варианты
строятся с помощью нелинейных функций относительных состояний.
